\section{The original arithmetic entries and their common marked frames}
\label{sec:arithmetic-frame-entries}

\paragraph{The endpoint and its retained premise.}
We prove the two original nineteen-field implications
\(\texttt{gen\_gd}\) and \(\texttt{gen\_host}\), then obtain the
original \(\texttt{GdFrameEx}\) and \(\texttt{HostFrameEx}\) from
\(P:\texttt{GC.Hypotheses}_{2}\).
The latter premise is exactly the cohomological package retained
in LFC22--23; it is not replaced by the assertion that a formal
proof exists.  The coefficient ring in these two entries is the
original coefficient DVR \(O=c.O\).
The Dat modulus is \(\ell N|\operatorname{disc}K|\), with
\(H_b=G_{\mathbb Q}\).  The Host modulus is \(\ell N\), with
\(H_b=\ker\operatorname{dir}(\kappa_B)\).
The Host argument allows \(\ell=3\), and both arguments work at
the valuation specified in the input.

The main simplifications are elementary.  Quadratic Gauss sums
place the imaginary quadratic field in the cyclotomic field of
its discriminant.  The two Dat nondegeneracy fields follow
directly from the original kernel and weight conditions.
A dyadic conductor calculation supplies the Host witness at
every relevant local place, without first changing the place
or assuming an auxiliary Host frame exists.  We also construct
that original auxiliary frame and prove its full
\(\texttt{exists\_sigma}\) consequence.

\paragraph{Notation and foundations.}
Let \(\Omega\) be the specified algebraic closure of \(\mathbb Q\),
and put \(G=\operatorname{Gal}(\Omega/\mathbb Q)\).
For a valuation subring \(B\) above a rational prime \(p\), write
\(D_B,I_B\) for its actual decomposition and inertia subgroups.
We use the usual extension-of-valuations dictionary: \(D_B\)
is the absolute Galois group of the local completion, \(I_B\)
is the kernel of its action on the residue algebraic closure,
and \(D_B/I_B=\widehat{\mathbb Z}\) with arithmetic Frobenius
acting as \(x\mapsto x^p\).
Finite and infinite Galois correspondence, the existence of
these prolongations, elementary integral-basis theory, and
Hensel's lemma for the complete ring \(\mathbb Z_\ell\) are
the classical foundations used below.  The cyclotomic,
character, discriminant and conductor calculations needed
from them are proved explicitly.

For a positive integer \(n\), let
\[
 U_n=(\mathbb Z/n\mathbb Z)^\times,\qquad
 c_n:G\longrightarrow U_n
\]
be the action on \(\mu_n\).  Let
\(\epsilon:G\to\mathbb Z_\ell^\times\) be the inverse-limit
cyclotomic character.  If \(O=c.O\), put
\[
 k=O/\mathfrak m_O,\qquad
 \omega:G\longrightarrow\mathbb F_\ell^\times
                         \longrightarrow k^\times .
\]
The character \(\omega\) is the source's literal
\(\texttt{omO}\).  For a Dirichlet character
\(\psi:U_N\to O^\times\), write
\[
 \operatorname{dir}(\psi)=\psi c_N,\qquad
 \Psi=\epsilon^e\operatorname{dir}(\psi),\qquad
 \overline\Psi=\operatorname{red}\circ\Psi.
                                                        \tag{AFE.1}
\]
The power of \(\epsilon\) in (AFE.1) uses the specified
\(\mathbb Z_\ell\)-algebra map to \(O\).
All finite characters are given the discrete finite-image
topology.  Their global pullbacks are locally constant.

\subsection{Residue fields, character components and cyclotomic inertia}

% BEGIN AFE1
\paragraph{AFE1. The coefficient residue field and torsion reduction.}
\label{claim:AFE1}
For the original coefficient \(c\), the field \(k\) is finite
of characteristic \(\ell\).  Reduction is injective on roots
of unity of order prime to \(\ell\), and kills roots of unity
of \(\ell\)-power order.  The reduction of the specified map
\(\mathbb Z_\ell\to O\) is the usual map
\(\mathbb Z_\ell\to\mathbb F_\ell\to k\).

\begin{proof}
The ring \(O\) is finite over \(\mathbb Z_\ell\) and is a
subring of \(\overline{\mathbb Q}_\ell\).
Every \(x\in O\) is therefore integral over \(\mathbb Z_\ell\).
If \(|x|_\ell>1\), a monic integral equation for \(x\) has
its leading term strictly larger in absolute value than
every other term, which is impossible.  Hence
\(|x|_\ell\le1\).  In particular \(\ell^{-1}\notin O\), so
\(\ell\) is a nonunit and lies in \(\mathfrak m_O\).
The quotient \(O/\ell O\) is a finite-dimensional
\(\mathbb F_\ell\)-space, and \(k\) is a quotient of it.
Thus \(k\) is finite and has characteristic \(\ell\).
For \(z\in\mathbb Z_\ell\), choose its integer residue
representative \(a\).  The difference \(z-a\) belongs to
\(\ell\mathbb Z_\ell\), so its image in \(k\) vanishes.
This proves the exact reduction convention.

More generally, let \(R\) be any local ring with residue
characteristic \(\ell\).  If \(x^n=1\), \(\ell\nmid n\),
and \(\bar x=1\), then
\[
 (x-1)(1+x+\cdots+x^{n-1})=0.
\]
The second factor reduces to \(n\ne0\) and is a unit.
Thus \(x=1\).  Applying this to the ratio of two roots
proves injectivity on prime-to-\(\ell\) torsion.
In a field of characteristic \(\ell\),
\(X^{\ell^a}-1=(X-1)^{\ell^a}\); this proves the
assertion about \(\ell\)-power roots.

We will also use the elementary fact that a finite subgroup
of the multiplicative group of a field is cyclic.  Indeed,
if \(d\) is its exponent, the primary decomposition of a
finite abelian group produces an element of order \(d\).
Every element is a root of \(X^d-1\), so the subgroup has
at most \(d\) elements and hence is generated by that
element.  In particular all the finite character orders
used below are positive.
\end{proof}
% END AFE1

% BEGIN AFE2
\paragraph{AFE2. The exact modulus and image-primary decompositions.}
\label{claim:AFE2}
Write \(N=\prod_p p^{a_p}\).
Let \(\operatorname{pr}_p:U_N\to U_N\) retain the
\(p^{a_p}\)-component in the Chinese remainder decomposition
and set every other component to \(1\).  Then
\[
 \psi_p=\psi\operatorname{pr}_p=\texttt{comp }p\,\psi,\qquad
 \psi=\prod_{p\mid N}\psi_p.                              \tag{AFE.2}
\]
The conductor of \(\psi_p\) is \(p^{b_p}\), with
\(0\le b_p\le a_p\).  In particular the dyadic character
\(\psi_2\) has \(2\)-power order and conductor
\(2^{c_2(\psi)}\).

The source's \(\kappa=\texttt{kap }\ell\,\psi\) and
\(\theta=\texttt{th0 }\ell\,\psi\) are exactly the
\(\ell\)-primary and prime-to-\(\ell\) parts of the character:
\[
 \psi=\theta\kappa,\quad
 \kappa_L=\kappa_\ell,\quad
 \kappa_B=\kappa\kappa_L^{-1},\quad
 \psi=\theta\kappa_B\kappa_L.                             \tag{AFE.3}
\]
For odd \(\ell\), \(\kappa_2=1\), and consequently
\(\psi_2=\theta_2\).

\begin{proof}
Put \(A=p^{a_p}\) and \(C=N/A\).  The integer
\(C^{\varphi(A)}\) is \(1\) modulo \(A\), and is \(0\)
modulo \(C\); when \(A=1\) it is \(0\) modulo \(N\).
Thus the source's idempotent
\(\texttt{idem }p\,N\) has exactly those two components.
The formula
\[
 u\longmapsto \texttt{idem }p\,N\cdot u+
                         (1-\texttt{idem }p\,N)
\]
is precisely \(\operatorname{pr}_p\).  The product of these
projections is \(u\), proving (AFE.2), as well as
idempotence, commutation of different projections, and
compatibility with products and inverses of characters.
A character depending only on the \(p\)-component factors
through the least quotient \(U_{p^{b_p}}\) on which it is
well defined.  This is its conductor: triviality on the
kernel is exactly the condition for factoring through a
surjective group map.  Its conductor therefore divides
\(p^{a_p}\).  The group \(U_{2^{a_2}}\) has \(2\)-power
order, proving the dyadic assertion.

For a finite-order element \(x\), write its order as
\(\ell^a d\), with \(\ell\nmid d\).
The exponent \(d^{\varphi(\ell^a)}\) is \(1\) modulo
\(\ell^a\) and is \(0\) modulo \(d\).  Thus
\[
 x^{\,d^{\varphi(\ell^a)}}
\]
is exactly the \(\ell\)-primary part of \(x\).
This is the literal \(\texttt{lPart}\) formula, including
the case \(a=0\).  Apply it in the finite character group
to obtain \(\kappa\); its complementary factor is
\(\theta=\psi\kappa^{-1}\).
Projection to modulus components commutes with taking
these powers, giving (AFE.3).
Finally \(\kappa_2\) has both \(\ell\)-power and
\(2\)-power order.  For odd \(\ell\) it is trivial.
\end{proof}
% END AFE2

% BEGIN AFE3
\paragraph{AFE3. Cyclotomic inertia and corrected Frobenius at every place.}
\label{claim:AFE3}
For every \(B\) above \(p\), the image of \(I_B\) under
\(c_N\) is precisely \(\operatorname{pr}_p(U_N)\).
If the integer \(m\) is a unit in a valuation subring,
its inertia fixes \(\mu_m\).
For every \(B\) above \(\ell\),
\(\epsilon(I_B)=\mathbb Z_\ell^\times\).
For each \(B\) above \(p\) there is \(\gamma\in D_B\)
whose action on every root of unity of order prime to
\(p\) is \(\zeta\mapsto\zeta^p\).  If \(p=\ell\), it
can be chosen with \(\epsilon(\gamma)=1\).
If \(p\ne\ell\) and \(p\equiv1\pmod\ell\), an ordinary
Frobenius lift already has \(\omega(\gamma)=1\).

\begin{proof}
A root of unity and its inverse are integral at every
valuation, so both have valuation \(1\).
If inertia sends a root \(\zeta\in\mu_m\) to \(\zeta'\),
their residues agree.  The ratio \(x=\zeta'/\zeta\)
satisfies \(x^m=1\) and has residue \(1\).
When \(m\) is a unit the geometric-sum argument in
AFE1, applied to this valuation ring, gives \(x=1\).
Thus inertia fixes all prime-to-\(p\) roots at \(B\).

For \(a\ge1\), the polynomial
\(\Phi_{p^a}(X+1)\) is Eisenstein at \(p\).
Its reduction is \(X^{p^{a-1}(p-1)}\), its constant
term is \(p\), and its leading coefficient is \(1\).
It follows that \(\mathbb Q_p(\mu_{p^a})/\mathbb Q_p\)
has degree \(\varphi(p^a)\) and is totally ramified.
The faithful cyclotomic action therefore identifies its
Galois group with \(U_{p^a}\), and inertia surjects onto
that group.  Together with the preceding prime-to-\(p\)
assertion and the Chinese remainder theorem, this proves
the exact image \(\operatorname{pr}_p(U_N)\).
The argument takes place in the completion belonging to
the given \(B\), so it applies to every prolongation.

The compatible maps \(I_B\to U_{\ell^a}\) are all onto.
For any specified element of \(\mathbb Z_\ell^\times\),
the sets of its lifts at levels \(a\) are nonempty nested
closed subsets of the compact group \(I_B\).
Their intersection is nonempty.  This proves
\(\epsilon(I_B)=\mathbb Z_\ell^\times\).

Choose a lift \(\gamma\in D_B\) of arithmetic Frobenius
in \(D_B/I_B=\widehat{\mathbb Z}\).
For a prime-to-\(p\) root \(\zeta\), the elements
\(\gamma\zeta\) and \(\zeta^p\) have the same residue;
injectivity of reduction on those roots gives equality.
This holds simultaneously for all such roots.
If \(p=\ell\), choose \(\tau\in I_B\) with
\(\epsilon(\tau)=\epsilon(\gamma)^{-1}\) and replace
\(\gamma\) by \(\gamma\tau\).  Inertia acts trivially
on the prime-to-\(p\) roots, so their Frobenius action
is unchanged and the new cyclotomic character is \(1\).
If \(p\ne\ell\), Frobenius acts on \(\mu_\ell\)
by the exponent \(p\), which proves the last assertion.
\end{proof}
% END AFE3

% BEGIN AFE4
\paragraph{AFE4. A quadratic discriminant square root with the exact sign character.}
\label{claim:AFE4}
Let \(K\) be the original imaginary quadratic field,
embedded in \(\Omega\), and let
\(\eta:G\to\{\pm1\}\) have kernel \(G_K\).
Write \(d=\operatorname{disc}K\) and \(D=|d|\).
Then \(D>0\), and there is \(\delta\in K\subset\Omega\)
such that
\[
       \delta^2=-D,\qquad s\delta=\eta(s)\delta
                          \quad(s\in G).                \tag{AFE.4}
\]

\begin{proof}
Write \(K=\mathbb Q(\sqrt m)\) for a negative squarefree
integer \(m\).  Here is the integral-basis calculation
including the possible dyadic factor.
If \(a+b\sqrt m\) is integral, its trace gives
\(2a\in\mathbb Z\), and its norm gives
\(4mb^2\in\mathbb Z\).
Squarefreeness of \(m\), prime by prime in the rational
denominator of \(b\), implies \(2b\in\mathbb Z\).
Set \(r=2a,s=2b\).  Integrality is then equivalent to
\[
                  r,s\in\mathbb Z,\qquad
                  r^2-ms^2\equiv0\pmod4.
\]
If \(m\equiv1\pmod4\), the two integers have the same
parity; otherwise both are even.  Conversely these
conditions make trace and norm integral and hence give
a monic quadratic equation.  It follows that
\[
 \mathcal O_K=
 \begin{cases}
  \mathbb Z[(1+\sqrt m)/2],&m\equiv1\pmod4,\\
  \mathbb Z[\sqrt m],&m\equiv2,3\pmod4.
 \end{cases}
\]
For the resulting integral basis \((1,u)\), the determinant
of the embedding matrix is \(u'-u\), where the prime denotes
the other embedding.  Thus
\[
        d=(u-u')^2=
        \begin{cases}m,&m\equiv1\pmod4,\\4m,&\text{otherwise}.
        \end{cases}
\]
In particular \(d<0\).  Take \(\delta=u-u'\).
It is nonzero and generates \(K\) over \(\mathbb Q\).
An element of \(G\) fixing \(K\) fixes \(\delta\);
one not fixing \(K\) exchanges its two roots
\(\delta,-\delta\).
These alternatives agree exactly with
\(\eta=1,-1\), since \(\ker\eta=G_K\).
This proves (AFE.4).
\end{proof}
% END AFE4

% BEGIN AFE5
\paragraph{AFE5. The discriminant cyclotomic containment and its inertia consequence.}
\label{claim:AFE5}
For the field and character in AFE4,
\[
                  K\subset\mathbb Q(\mu_D).
                                                        \tag{AFE.5}
\]
Consequently an element fixing \(\mu_D\) lies in
\(\ker\eta\), and if \(\ell\nmid D\) then
\(I_B\subset\ker\eta\) for every \(B\) above \(\ell\).

\begin{proof}
We prove the required quadratic containment directly.
For an odd prime \(p\), choose a primitive \(p\)-th
root \(\zeta\), let \(\chi\) be the quadratic character
of \(\mathbb F_p^\times\), extended by \(0\) at \(0\),
and put
\[
       \tau_p=\sum_{a\in\mathbb F_p}\chi(a)\zeta^a,\qquad
       \tau_p^-=\sum_a\chi(a)\zeta^{-a}.
 \]
The cyclicity in AFE1 gives
\(\chi(-1)=(-1)^{(p-1)/2}\).  Changing \(a\) to \(-a\)
shows \(\tau_p^-=\chi(-1)\tau_p\).
In the product \(\tau_p\tau_p^-\), change variables
\(a=tb\) for \(b\ne0\).  The inner sum
\(\sum_{b\ne0}\zeta^{(t-1)b}\) is \(p-1\) for
\(t=1\) and \(-1\) otherwise.  Since
\(\sum_t\chi(t)=0\), this gives
\[
              \tau_p\tau_p^-=p,\qquad
              \tau_p^2=(-1)^{(p-1)/2}p=:p^*.
                                                        \tag{AFE.6}
\]
In particular these roots are nonzero.

For odd squarefree \(m<0\), the product of the \(p^*\)
over \(p\mid |m|\) has absolute value \(|m|\) and is
congruent to \(1\) modulo \(4\).  It equals \(m\) if
\(m\equiv1\pmod4\), and \(-m\) if \(m\equiv3\pmod4\).
The product of the corresponding \(\tau_p\)'s therefore
gives \(\sqrt m\) inside \(\mathbb Q(\mu_{|m|})\) in
the first case; in the second, multiplying by \(i\)
gives it inside \(\mathbb Q(\mu_{4|m|})\).
The case \(m=-1\) uses the empty product and the same
factor \(i\).
For even squarefree \(m=2u<0\), first obtain
\(\sqrt u\) inside \(\mathbb Q(\mu_{4|u|})\) by the
odd calculation, and multiply by
\(\zeta_8+\zeta_8^{-1}\), whose square is \(2\).
This puts \(\sqrt m\) inside
\(\mathbb Q(\mu_{8|u|})\).
By the discriminant calculation in AFE4, the modulus
in each case is exactly \(D\).
This proves (AFE.5), including all dyadic cases.

An automorphism fixing \(\mu_D\) therefore fixes \(K\).
When \(\ell\nmid D\), AFE3 says that every inertia
group above \(\ell\) fixes \(\mu_D\).
The equality \(\ker\eta=G_K\) proves both consequences.
\end{proof}
% END AFE5

% BEGIN AFE6
\paragraph{AFE6. One inertia exponent, and the two Dat nondegeneracy fields.}
\label{claim:AFE6}
For every \(\psi\) there is \(a\ge0\), chosen independently
of the valuation \(B\), such that
\[
       \overline\Psi|_{I_B}=\omega^{e+a}|_{I_B}
                       \quad(B\mid\ell).                 \tag{AFE.7}
\]
If \(\ell\nmid\operatorname{ord}(\psi)\), the full
\(\ell\)-modulus component has the exact form
\[
       \operatorname{dir}(\psi_\ell)=\widetilde\omega^a,
                                                        \tag{AFE.8}
\]
where \(\widetilde\omega\) is the Teichmuller lift of
\(\omega\) through \(\mathbb Z_\ell^\times\).
For the original Dat input, these facts imply
\[
 \ker\omega\cap\ker\overline\Psi\not\subset\ker\eta,\qquad
 \forall B\mid\ell\ \exists s\in I_B:
                  \eta(s)=1,\quad\overline\Psi(s)^2\ne1.
                                                        \tag{AFE.9}
\]

\begin{proof}
By AFE3, restriction to \(I_B\) sees only \(\psi_\ell\).
The kernel of \(U_{\ell^b}\to\mathbb F_\ell^\times\)
is an \(\ell\)-group, and has trivial image in \(k^\times\)
by AFE1.  Thus the reduced character factors through
\(\mathbb F_\ell^\times\).  Every root of
\(X^{\ell-1}-1\) in \(k\) is in its prime field:
the \(\ell-1\) nonzero prime-field elements already
exhaust the roots.  The cyclicity of
\(\mathbb F_\ell^\times\) consequently writes this
character as \(x\mapsto x^a\).
This is a statement about the original component
\(\psi_\ell\), so one exponent works for all \(B\).
The exact cyclotomic reduction in AFE1 proves (AFE.7).

If the order of \(\psi\) is prime to \(\ell\), the wild
kernel also has trivial image before reduction:
its image has both \(\ell\)-power and prime-to-\(\ell\)
order.  Hence \(\psi_\ell\) itself factors through
\(\mathbb F_\ell^\times\).
For each \(u\in\mathbb F_\ell^\times\), Hensel's lemma
applied to \(X^{\ell-1}-1\) gives a unique lift
\([u]\in\mathbb Z_\ell^\times\) of order dividing
\(\ell-1\).  Uniqueness proves \([uv]=[u][v]\).
The roots \(\psi_\ell(u)\) and \([u]^a\) have the same
residue and order prime to \(\ell\), so AFE1 makes
them equal.  This proves (AFE.8).

Now retain the original Dat and all its GD fields.
The order hypothesis gives \(\theta=\psi\), so its
character \(\chi_P\) is precisely
\(\psi\psi_\ell^{-1}\).
The GD kernel hypothesis supplies \(t\in G\) with
\[
      c_\ell(t)=1,\qquad
      \operatorname{dir}(\chi_P)(t)=1,\qquad
      t\notin G_K.
 \]
By (AFE.8), \(\operatorname{dir}(\psi_\ell)(t)=1\).
Thus \(\operatorname{dir}(\psi)(t)=1\) and
\(\overline\Psi(t)=1\), while
\(\eta(t)\ne1\).  This proves the first assertion
of (AFE.9), with the literal \(\omega\)-kernel since
\(\mathbb F_\ell\to k\) is injective.

The GD condition \(\psi_\ell^2=1\), together with
surjectivity of \(I_B\to\mathbb F_\ell^\times\), gives
\(\ell-1\mid2a\).
Its weight condition is \(\ell-1\nmid2e\); hence
\(\ell-1\nmid2(e+a)\).
Choose \(s\in I_B\) whose \(\omega(s)\) generates
\(\mathbb F_\ell^\times\).  Equation (AFE.7) gives
\(\overline\Psi(s)^2\ne1\).
The original coprimality
\(\ell\nmid D\,h_K\) implies \(\ell\nmid D\),
so AFE5 gives \(\eta(s)=1\).
This proves the second assertion at every \(B\).
Only the Dat branch uses this weight condition;
it imposes no restriction on the Host prime.
\end{proof}
% END AFE6

\subsection{The full Dat entry}

% BEGIN AFE7
\paragraph{AFE7. All nineteen original fields of \(\texttt{gen\_gd}\).}
\label{claim:AFE7}
Let the input be exactly the original
\(\texttt{Dat}(\Omega,K,c,e,\psi)\), with its specified
embedding of the imaginary quadratic \(K\), and let
\(B\) be any valuation above \(\ell\).
Then
\[
 \texttt{Gen}\bigl(\ell,\eta,\Psi,\mathrm{cc},B,
                     \ell N|\operatorname{disc}K|,G\bigr)
                                                        \tag{AFE.10}
\]
holds.  In particular this is the full original
\(\texttt{gen\_gd}\), with \(H_b=\top\).

\begin{proof}
Put \(D=|\operatorname{disc}K|\) and \(M=\ell ND\).
We verify the actual nineteen fields.

\emph{The fields \(\texttt{odd},\texttt{fin},
\texttt{M\_ne},\texttt{l\_dvd},\texttt{two\_dvd},
\texttt{cc},\texttt{over\_l}\).}
The prime is odd by GD.  AFE1 gives finiteness of \(k\).
The integers \(\ell,N,D\) are positive, so \(M\ne0\)
and \(\ell\mid M\).  GD says \(N\) is even, hence
\(2\mid M\).  The Dat embedding into \(\mathbb C\)
realizes its specified \(\mathrm{cc}\) as complex
conjugation.  The valuation condition is the given one.

\emph{The field \(\texttt{quad}\).}
Use \(D,\delta\) from AFE4.  Every prime dividing
\(D\) divides \(M\), and (AFE.4) is the exact
required square and transformation law.

\emph{The field \(\texttt{cyc}\).}
An element fixing \(\mu_M\) fixes \(\mu_D\), and
therefore has \(\eta=1\) by AFE5.
It also fixes \(\mu_N\) and \(\mu_\ell\), so
\(\operatorname{dir}(\psi)=1\) and \(\omega=1\).
By (AFE.1) and AFE1, \(\overline\Psi=1\).
Thus the required two \(\texttt{KillsMu}\) clauses
hold with exponent \(1\).

\emph{The fields \(\eta_\ell,\texttt{even},\Psi_\ell\).}
The first is AFE5 using \(\ell\nmid D\).
Complex conjugation acts by inversion on every root
of unity, so
\[
       \epsilon(\mathrm{cc})=-1,\qquad
       \operatorname{dir}(\psi)(\mathrm{cc})=\psi(-1).
 \]
The original parity condition
\(\psi(-1)=(-1)^e\) gives \(\Psi(\mathrm{cc})=1\)
even before reduction.  The common exponent
\(j=e+a\) in AFE6 is precisely \(\Psi_\ell\).

\emph{The fields \(\texttt{g4},\texttt{g5}\).}
These are (AFE.9).  Its second witness lies in
\(I_B\subset D_B\), so it satisfies the decomposition
condition at the specified \(B\).

\emph{The four \(H_b\) fields.}
The subgroup \(G\) is open, contains the commutator
and every inertia group, and has index
\(1=\ell^0\).  These give
\(\texttt{Hb\_open},\texttt{Hb\_ab},\texttt{Hb\_idx},
\texttt{Hb\_l}\) literally.

\emph{The field \(\texttt{teich}\).}
Let \(B'\) be any valuation not above \(\ell\)
and \(s\in I_{B'}\).  The integer \(\ell\) is
a unit in this valuation ring, so AFE3 shows
that \(s\) fixes every \(\ell\)-power root.
Thus \(\epsilon(s)=1\) and
\(\Psi(s)=\operatorname{dir}(\psi)(s)\).
The latter has order prime to \(\ell\) by GD.
If its reduction is \(1\), AFE1 makes it \(1\).
This also covers a trivial valuation, and proves
the exact quantification over all \(B'\).
All nineteen fields have now been supplied.
\end{proof}
% END AFE7

\subsection{The dyadic Host arguments}

% BEGIN AFE8
\paragraph{AFE8. The dyadic conductor forces the Host kernel obstruction.}
\label{claim:AFE8}
Suppose \(\ell\) is odd and \(c_2(\psi)\ge4\).
There exists \(s\in G\) such that
\[
       \omega(s)=1,\qquad \overline\Psi(s)=1,\qquad
                         \eta_4(s)=-1,                 \tag{AFE.11}
\]
where \(\eta_4\) is the quadratic character of
\(\mathbb Q(i)\).  This proves the full Host
\(\texttt{g4}\) obstruction.

\begin{proof}
The conductor of \(\psi_2\) is \(2^{c_2(\psi)}\)
by AFE2, so \(4\mid N\).
If \(\psi_2\) were trivial on the kernel of
\(U_N\to U_4\), it would factor through that
surjective map and have conductor dividing \(4\).
Therefore there is \(v\in U_N\), \(v\equiv1\pmod4\),
with \(\psi_2(v)\ne1\).
This value has order \(2^{b+1}\) for some \(b\ge0\).
For \(w=v^{2^b}\) we have \(w\equiv1\pmod4\)
and \(\psi_2(w)=-1\): in the characteristic-zero
domain \(O\), the only nonidentity element whose
square is \(1\) is \(-1\).
The value \(\psi_2(-1)\) is either \(1\) or \(-1\).
In the first case put \(u=-1\), and in the second
put \(u=-w\).  In both cases
\[
                    u\equiv-1\pmod4,\qquad
                    \psi_2(u)=1.                       \tag{AFE.12}
\]

Choose any prolongation of the \(2\)-adic valuation.
AFE3 supplies \(s\) in its inertia group with
\(c_N(s)=\operatorname{pr}_2(u)\).
As \(\ell\) is odd, the same inertia fixes
\(\mu_\ell\), so \(\omega(s)=1\).
By (AFE.12),
\(\operatorname{dir}(\psi)(s)=\psi_2(u)=1\),
whence \(\overline\Psi(s)=1\).
Its action on \(\mu_4\) has exponent \(-1\),
so it sends \(i\) to \(-i\) and has
\(\eta_4(s)=-1\).
This proves (AFE.11) for every exponent \(e\).
\end{proof}
% END AFE8

% BEGIN AFE9
\paragraph{AFE9. The dyadic power bound from the actual conductor.}
\label{claim:AFE9}
Let \(p>1\) be odd, \(o\ge1\), and suppose
\[
                  8op^2<2^{c_2(\psi)}.                 \tag{AFE.13}
\]
If the dyadic component of \(u\in U_N\) is the
residue of \(p\), then
\[
                         \psi_2(u)^{4o}\ne1.           \tag{AFE.14}
\]
The argument only requires that the character values
lie in a characteristic-zero domain.

\begin{proof}
We first give the principal-unit calculation.
If \(t\ge2\) and \(b\) is odd, repeated squaring gives
\[
             (1+2^t b)^{2^h}
                    =1+2^{t+h}b_h,\qquad b_h\ \text{odd}.
                                                        \tag{AFE.15}
\]
For the induction step, the new odd coefficient is
\(b_h+2^{t+h-1}b_h^2\).
It follows that, modulo \(2^A\) with \(A>t\),
the element \(1+2^t b\) has exact order \(2^{A-t}\).
Indeed (AFE.15) gives this for powers of \(2\);
an odd additional exponent preserves the valuation
of the difference from \(1\), by the geometric sum.
The group of units congruent to \(1\) modulo \(2^t\)
has exactly \(2^{A-t}\) elements.  Thus this element
generates that entire group.

Write \(o=2^j m\) with \(m\) odd, and
\(p^2-1=2^s b\) with \(b\) odd.
Since \(p\) is odd, \(s\ge3\).
Assume, for contradiction, that (AFE.14) fails.
The value \(\psi_2(u)\) has \(2\)-power order;
cancelling the odd factor \(m\) in its order gives
\[
                         \psi_2(u)^{2^{j+2}}=1.
 \]
By (AFE.15), the dyadic component of
\(u^{2^{j+2}}\) is
\[
         p^{2^{j+2}}=(p^2)^{2^{j+1}}
                    =1+2^{s+j+1}b',\qquad b'\ \text{odd}.
 \]
Put \(t=s+j+1\) and \(c=c_2(\psi)\).
The numerical estimate
\[
        2^t=2\cdot2^j\,2^s
                    <2op^2\le8op^2<2^c
 \]
shows \(t<c\).
If \(2^A\) is the exact dyadic divisor of \(N\),
then \(c\le A\) by AFE2.
The principal-unit calculation shows that the dyadic
component of \(u^{2^{j+2}}\) generates
\(\ker(U_{2^A}\to U_{2^t})\).
Our assumed character equality says that \(\psi_2\)
is trivial on this generator, hence on that kernel.
Since \(\psi_2\) is trivial on all other modulus
components, it factors through \(U_{2^t}\).
Its conductor therefore divides \(2^t\), contrary
to \(c>t\).
This proves (AFE.14) with the actual conductor,
not merely an upper bound on it.
\end{proof}
% END AFE9

% BEGIN AFE10
\paragraph{AFE10. A nonquadratic decomposition witness at every Host place.}
\label{claim:AFE10}
Retain the original
\(\texttt{HostData}(\ell,e,\psi,o)\).
For every rational prime \(p\in T_b(\ell,\psi)\)
and every valuation \(B\) above \(p\), there exists
\(x\in D_B\) with
\[
 \eta_4(x)=1,\qquad c_\ell(x)=1,\qquad
 \overline\Psi(x)^2\ne1,\qquad
                         \Psi(x)^2-1\in O^\times.
                                                        \tag{AFE.16}
\]
There is no restriction \(\ell>3\).

\begin{proof}
We first spell out which primes occur.
The set \(T_b\) is \(\{\ell\}\cup S_b\), where
\[
 S_b=\{p\mid N:p\ne\ell,\ \kappa_p\ne1\}.
 \]
For \(p\in S_b\), the nontrivial character
\(\kappa_p\) has \(\ell\)-power order, so
\(\ell\mid |U_{p^{a_p}}|=p^{a_p-1}(p-1)\).
Since \(p\ne\ell\), this gives
\(\ell\mid p-1\).  In particular \(p\) is odd and
\(p\equiv1\pmod\ell\).
The other possibility is \(p=\ell\), which is
also odd by HostData.

By AFE2, \(\psi_2=\theta_2\).
The HostData order field states that each odd
modulus component \(\theta_r\) has order dividing
\(2o\).  Applying (AFE.2) gives the character identity
\[
                         \theta^{4o}=\psi_2^{4o}.
                                                        \tag{AFE.17}
\]
Choose \(\gamma\in D_B\) by AFE3.  If \(p=\ell\),
use the corrected Frobenius with \(\epsilon(\gamma)=1\);
if \(p\ne\ell\), use \(p\equiv1\pmod\ell\).
In both cases \(c_\ell(\gamma)=1\), and for
\(u=c_N(\gamma)\) the dyadic component is \(p\).
The field \(\texttt{R}\) in HostData supplies (AFE.13).

Suppose that \(\overline{\psi(u)}^{\,4}=1\).
The character \(\kappa\) reduces to \(1\) by AFE1,
so \(\overline{\theta(u)^4}=1\).
The value \(\theta(u)^4\) has order prime to \(\ell\);
AFE1 gives \(\theta(u)^4=1\).
Equation (AFE.17) would then give
\(\psi_2(u)^{4o}=1\), contradicting AFE9.
Consequently
\[
             \overline{\psi(u)}^{\,4}\ne1.
 \]
Set \(x=\gamma^2\).
It still belongs to \(D_B\);
the sign character gives \(\eta_4(x)=1\), and
\(c_\ell(x)=1\).
Using \(\omega(\gamma)=1\) in (AFE.1) yields
\[
            \overline\Psi(x)^2
                     =\overline{\psi(u)}^{\,4}\ne1.
 \]
In a local ring an element is a unit precisely when
its residue is nonzero.  Therefore
\(\Psi(x)^2-1\) is a unit.  The construction was
performed in the specified \(D_B\); no change of
valuation or globally chosen substitute was made.
\end{proof}
% END AFE10

% BEGIN AFE11
\paragraph{AFE11. The literal Host subgroup and the away-from-\(\ell\) implication.}
\label{claim:AFE11}
For odd \(\ell\), define
\[
                   H_b=\ker\operatorname{dir}(\kappa_B).
                                                        \tag{AFE.18}
\]
Then \(H_b\) is open, contains the commutator of \(G\),
has index a power of \(\ell\), and contains every
inertia group above \(\ell\).
For every valuation \(B'\) not above \(\ell\),
\[
 s\in I_{B'}\cap H_b,\quad\overline\Psi(s)=1
                         \ \Longrightarrow\ \Psi(s)=1.
                                                        \tag{AFE.19}
\]

\begin{proof}
The character \(\kappa_B\) has finite \(\ell\)-power
order and factors through \(U_N\).
Its global pullback is locally constant with finite
abelian \(\ell\)-power image.
Its kernel is therefore open, contains the commutator,
and has index equal to the cardinality of that image,
which is \(\ell^a\) for some \(a\ge0\).
The \(\ell\)-modulus component of \(\kappa_B\) is
\(\kappa_\ell\kappa_L^{-1}=1\).
By AFE3 its pullback is trivial on every \(I_B\)
above \(\ell\), proving the inertia containment.

For \(s\in I_{B'}\) as in (AFE.19), AFE3 gives
\(\epsilon(s)=1\) and
\(\operatorname{dir}(\kappa_L)(s)=1\):
the first follows from all \(\ell\)-power roots,
and the second depends on only a finite level of them.
Membership in \(H_b\) kills \(\kappa_B\).
Equation (AFE.3) therefore gives
\[
                         \Psi(s)=\operatorname{dir}(\theta)(s).
 \]
The right side has order prime to \(\ell\).
Its residue being \(1\), AFE1 makes it \(1\).
The argument uses only that \(\ell\) is a unit in
the valuation ring, so it covers every \(B'\) in
the original field's quantifier.
\end{proof}
% END AFE11

% BEGIN AFE12
\paragraph{AFE12. All nineteen original fields of \(\texttt{gen\_host}\).}
\label{claim:AFE12}
For the original HostData, any specified complex
conjugation \(\sigma\), and any specified \(B\) above
\(\ell\), the subgroup (AFE.18) satisfies
\[
        \texttt{Gen}(\ell,\eta_4,\Psi,\sigma,B,\ell N,H_b).
                                                        \tag{AFE.20}
\]
This proves the full original existential
\(\texttt{gen\_host}\), including \(\ell=3\).

\begin{proof}
Set \(M=\ell N\) and keep precisely the subgroup
(AFE.18).  We again check all nineteen fields.

\emph{The first seven fields.}
HostData gives \(\ell\ne2\) and even \(N\).
AFE1 gives finite \(k\).
As \(\ell\) is prime and \(N>0\),
\(M\ne0\), \(\ell\mid M\), and \(2\mid M\).
The supplied \(\sigma\) and \(B\) give
\(\texttt{cc}\) and \(\texttt{over\_l}\).

\emph{The field \(\texttt{quad}\).}
Choose \(i\in\Omega\) with \(i^2=-1\), and use the
positive integer \(D=1\).
The prime-divisor condition is vacuous.
The action on \(\mu_4\) is \(i\mapsto i\) or
\(i\mapsto-i\); by the definition of the character
\(\chi_4:U_4\to\{\pm1\}\), this is
\(s(i)=\eta_4(s)i\).

\emph{The field \(\texttt{cyc}\).}
Use exponent \(2\).
Since \(M\) is even, \(4\mid M^2\);
also \(N\mid M^2\) and \(\ell\mid M^2\).
Fixing \(\mu_{M^2}\) therefore kills \(\eta_4\),
\(\operatorname{dir}(\psi)\), and \(\omega\).
Equation (AFE.1) gives the required two
\(\texttt{KillsMu}\) conclusions.

\emph{The fields \(\eta_\ell,\texttt{even},\Psi_\ell\).}
As \(\ell\) is odd, AFE3 says that inertia at
\(\ell\) fixes \(\mu_4\), proving \(\eta_\ell\).
Complex conjugation acts by inversion on roots;
the original parity condition again gives
\(\Psi(\sigma)=(-1)^e\psi(-1)=1\).
AFE6 supplies a single \(j=e+a\) before all
valuations, which is exactly \(\Psi_\ell\).

\emph{The fields \(\texttt{g4},\texttt{g5}\).}
HostData's \(c_2(\psi)\ge4\) and AFE8 prove
\(\texttt{g4}\).
The prime \(\ell\) belongs to \(T_b\), so AFE10
at the specified \(B\) proves \(\texttt{g5}\).

\emph{The four \(H_b\) fields and \(\texttt{teich}\).}
They are precisely AFE11.
Thus no \(H_b\) condition has been replaced by the
corresponding easier condition for \(\top\).
The conductor and Frobenius arguments require only
that \(\ell\) be odd, so all these fields apply when
\(\ell=3\).
\end{proof}
% END AFE12

\subsection{The original auxiliary Host frame and its local witnesses}

% BEGIN AFE13
\paragraph{AFE13. Constructing the full original \(\texttt{IsHostFrame}\).}
\label{claim:AFE13}
Under the original HostData there exist the finite
set \(S\), actual prolongations \(A_v\), and the
original frame \(F\) satisfying all ten
\(\texttt{IsHostFrame}\) fields.
This is the full \(\texttt{exists\_hostFrame}\);
it does not require \(P\).

\begin{proof}
Identify finite places of \(\mathbb Q\) with rational
prime ideals, and let
\[
       S=\{v:v\mid\ell N\}.
 \]
This is finite because \(\ell N\ne0\).
Choose a prolongation \(A_v\) of each \(v\in S\).
Let \(R_S\) be the closed normal subgroup of \(G\)
generated by all inertia groups at places outside
\(S\), and set
\[
       \Gamma=G/R_S=G_{\mathbb Q,S},\qquad
       \pi:G\twoheadrightarrow\Gamma.
 \]
This is exactly the source's \(\texttt{GS}\) and
\(\texttt{toGS}\), not a different quotient.
An inertia group outside \(S\) fixes \(\mu_N\),
\(\mu_4\), and every \(\ell\)-power root by AFE3:
for \(\mu_4\), use \(2\mid N\).
The kernels of \(c_N,c_4\) are open and closed,
and the kernel of \(\epsilon\) is an intersection
of such closed kernels.  Normality and closedness
therefore show that all three characters kill \(R_S\).
They descend to characters
\[
       \mu_N:\Gamma\to U_N,\quad
       \mu_4:\Gamma\to U_4,\quad
       \epsilon_S:\Gamma\to\mathbb Z_\ell^\times
 \]
whose pullbacks along \(\pi\) are the original ones.

Choose an embedding \(\Omega\hookrightarrow\mathbb C\);
its image is the field of algebraic numbers and is
stable under complex conjugation.
Conjugation therefore pulls back to
\(\sigma\in G\) with \(\texttt{IsCC}\,\sigma\).
Define the actual frame by
\[
 \begin{gathered}
 F.\pi=\pi,\qquad F.\eta=\chi_4\mu_4,\qquad
 F.\Psi=\epsilon_S^e\,\psi\mu_N,\qquad F.c=\pi(\sigma),\\
 F.D_v=D_{A_v},\qquad F.I_v=I_{A_v},\qquad
 F.v_\ell=(\ell),\\
 F.S_b=\{v\in S:\exists p\in S_b(\ell,\psi),\ v=(p)\}.
 \end{gathered}                                         \tag{AFE.21}
\]
The distinguished place is in \(S\), and each prime
in \(S_b(\ell,\psi)\) divides \(N\), so these choices
have the required types.
In the source's frame the groups \(F.D_v,F.I_v\)
are subgroups of the original \(G\); their maps to
\(\Gamma\) are induced by \(\pi\), just as in (AFE.21).

The ten fields can now be read off without additional
choices: membership of \(S\) is exactly divisibility
by \(\ell N\); \(A_v\) is over \(v\);
the quotient map is literally \(\texttt{toGS}\);
the two character pullbacks are \(\eta_4,\Psi\);
\(\sigma\) gives the complex-conjugation field;
the decomposition and inertia groups are the actual
ones; \(\ell\) belongs to the distinguished prime;
and the displayed \(S_b\) is its exact required
prime-divisibility condition.
These are respectively
\(\texttt{mem},\texttt{isOver},\pi,\eta,\Psi,
\texttt{cc},D,I,\texttt{vl},\texttt{Sb}\).
\end{proof}
% END AFE13

% BEGIN AFE14
\paragraph{AFE14. The full original \(\texttt{exists\_sigma}\) on any Host frame.}
\label{claim:AFE14}
Let \(F,S,(A_v)\) be any supplied original
\(\texttt{IsHostFrame}\), not necessarily the one
chosen in AFE13.
For every \(v\in F.T_b=\{F.v_\ell\}\cup F.S_b\)
there is \(x\in F.D_v\) with
\[
          F.\eta(F.\pi x)=1,\qquad
          F.\Psi(F.\pi x)^2-1\in O^\times,\qquad
                              c_\ell(x)=1.              \tag{AFE.22}
\]

\begin{proof}
If \(v=F.v_\ell\), the field \(\texttt{vl}\)
identifies its rational prime with \(\ell\).
If \(v\in F.S_b\), the field \(\texttt{Sb}\)
supplies \(p\in S_b(\ell,\psi)\) with \(v=(p)\).
In either case \(v=(p)\) for \(p\in T_b(\ell,\psi)\).
The field \(\texttt{isOver}\) puts \(A_v\) above
that exact \(p\), so AFE10 gives \(x\in D_{A_v}\)
with (AFE.16).
Use the literal equalities
\[
        F.D_v=D_{A_v},\quad
        F.\eta\circ F.\pi=\eta_4,\quad
        F.\Psi\circ F.\pi=\Psi
 \]
to obtain all three assertions (AFE.22).
AFE1 supplies the residue characteristic even when
it is not separately listed in this source theorem.
Thus this argument proves the full supplied-frame
statement and is compatible with every allowed choice
of the local prolongations.
\end{proof}
% END AFE14

\subsection{Joining the entries to the completed common frame}

% BEGIN AFE15
\paragraph{AFE15. The full original \(\boldsymbol{P\Rightarrow\texttt{GdFrameEx}}\).}
\label{claim:AFE15}
Assume the original \(P:\texttt{GC.Hypotheses}_{2}\).
For every original Dat input and every specified
\(B\) above \(\ell\), there are one CM field \(L\),
one \(g_1\in G\), one frame \(\Phi\), and one \(C_0\)
such that, with \(M=\ell N|\operatorname{disc}K|\),
\[
 \operatorname{GalLink}(\eta,\Psi,\mathrm{cc},g_1,B,M,L,\Phi),
 \quad \Phi.\texttt{Ok}(C_0),\quad \Phi.\texttt{Big}(C_0),
 \quad \texttt{SeedClause}(c,\eta,\Psi,\mathrm{cc},L).
                                                        \tag{AFE.23}
\]
This is the original \(\texttt{HasFrame}\) conclusion
and hence the full \(\texttt{GdFrameEx}\).

\begin{proof}
AFE7 supplies the complete original Gen with \(H_b=\top\).
Apply the inherited human character construction CAF7
to that input first.  It chooses a positive modulus
\(Q\) and a seed-producing function before choosing
the eventual field \(L\) and auxiliary prime \(w\).
Apply LFC22 with this \(Q\), the same Gen and the
original \(P\).  It supplies one
\[
               L,\ g_1,\ \Phi,\ C_0,\ w
 \]
with actual Linked data, all twenty-eight Ok fields,
Big for the same \(C_0\), the original Frobenius
witness at \(w\), \(w\equiv1\pmod Q\), and
\(\operatorname{upl}(L)\subset H^{\mathrm{th}}(w)\).
Let \(k_0\) denote its GalLink.
Its abelian and fixed-field clauses and those
conditions at \(w\) are exactly the field conditions
required by CAF7's previously chosen function.

We verify the remaining marked-element condition
on this very \(L,g_1\).
The GalLink field \(\texttt{g0}\) gives
\(g_1\in\operatorname{upl}(L)\) and
\(k_0.p(g_1)=\Phi.g_0\).
Its sign relation and \(\Phi.g_0\in\ker\Phi.\delta\)
give \(\eta(g_1)=1\).
The fixed-field equality in GalLink consequently
gives \(g_1\in G_L\).
Its character relation gives
\[
             \overline\Psi(g_1)=\Phi.\Psi(\Phi.g_0).
 \]
The actual Ok field \(\texttt{baseSubgroup\_spec}\)
says that the square of the right side is not \(1\).
Thus \(g_1\) supplies exactly the nonquadratic marker
needed by CAF11 to install the seed.
Apply CAF11, equivalently the final installation in
LFC23, to the already chosen \(L,w\).
It returns the actual SeedClause on this same \(L\);
it chooses neither a new field nor a new marker.
Together with the GalLink, Ok and Big already
obtained, this is (AFE.23).
Every quantified original Dat and every specified
\(B\) were arbitrary, so the full \(\texttt{GdFrameEx}\)
follows.
\end{proof}
% END AFE15

% BEGIN AFE16
\paragraph{AFE16. The full original \(\boldsymbol{P\Rightarrow\texttt{HostFrameEx}}\).}
\label{claim:AFE16}
Assume the same original \(P\).
For every original HostData, every specified complex
conjugation \(\sigma\), and every specified \(B\)
above \(\ell\), there are one CM field \(L\), one
\(g_1\), one \(\Phi\), and one \(C_0\) satisfying
\[
 \operatorname{GalLink}(\eta_4,\Psi,\sigma,g_1,B,\ell N,L,\Phi),
 \quad\Phi.\texttt{Ok}(C_0),\quad\Phi.\texttt{Big}(C_0),
 \quad\texttt{SeedClause}(c,\eta_4,\Psi,\sigma,L).
                                                        \tag{AFE.24}
\]
This is the full original \(\texttt{HostFrameEx}\),
with \(\ell=3\) included.

\begin{proof}
AFE12 supplies full Gen with the actual subgroup
\(H_b=\ker\operatorname{dir}(\kappa_B)\) and
the exact modulus \(M=\ell N\).
Use CAF7 to choose its \(Q\) and character data
before the field and prime.
LFC22, at this \(Q\) and this Gen, supplies
the actual linked \(L,g_1,\Phi,C_0,w\), all Ok
fields, Big, the Frobenius witness,
\(w\equiv1\pmod Q\), and the required subgroup
inclusion.  The character construction is being
used at the original coefficient DVR \(O=c.O\);
we do not extend its seed assertion to a general
local coefficient ring.

For completeness the marking check is unchanged
in content but uses the Host character:
\(\texttt{g0}\) gives \(g_1\in\operatorname{upl}(L)\)
and \(k_0.p(g_1)=\Phi.g_0\);
\(\texttt{del}\) gives \(\eta_4(g_1)=1\);
\(\texttt{gal}\) gives \(g_1\in G_L\);
and \(\texttt{psi}\), followed by
\(\texttt{baseSubgroup\_spec}\), gives
\(\overline\Psi(g_1)^2\ne1\).
CAF11 therefore installs its previously selected
character seed on this same \(L,w\).
Keeping the GalLink, Ok and Big from the same
construction proves (AFE.24).
No step invokes the Dat weight condition or
requires \(\ell>3\).
\end{proof}
% END AFE16

\subsection{The original marking and the bounded finite-point transport}

The remaining claims make a further, separate advance.
The representation-theoretic input \(\texttt{FinPt}\)
is retained, but its transport to \(\texttt{FinEnd}\)
is proved, including the actual normalized representation
and its local characters.  In this subsection \(O\) may
be any commutative local ring of residue characteristic
\(\ell\), as in the original transport theorem.
No coefficient-DVR hypothesis is inserted into that
general statement.

% BEGIN AFE17
\paragraph{AFE17. The same \(g_1\) satisfies all original marking fields.}
\label{claim:AFE17}
The \(g_1\) in AFE15 satisfies
\(\texttt{Marks}(\eta,\Psi,\mathrm{cc},g_1)\).
The \(g_1\) in AFE16 satisfies
\(\texttt{Marks}(\eta_4,\Psi,\sigma,g_1)\).
Thus the frame and its marker already have the
original downstream marking, not merely a residual
version of it.

\begin{proof}
The five fields of Marks are
\[
 \sigma^2=1,\qquad \eta(\sigma)=-1,\qquad
 \Psi(\sigma)=1,\qquad \eta(g_1)=1,\qquad
                         \Psi(g_1)^2-1\in O^\times.
                                                        \tag{AFE.25}
\]
Complex conjugation squares to the identity on
\(\Omega\).
In the Dat case it acts nontrivially on the imaginary
quadratic field, so its sign under \(\eta\) is \(-1\).
In the Host case it sends \(i\) to \(-i\), so its
sign under \(\eta_4\) is \(-1\).
The original parity condition gives
\(\Psi(\sigma)=1\) before reduction, as proved in
AFE7 and AFE12.
The common-marker calculation in AFE15 and AFE16
already gives
\(\eta(g_1)=1\) and \(\overline\Psi(g_1)^2\ne1\).
The latter says that \(\Psi(g_1)^2-1\) has nonzero
residue, hence is a unit of the local ring.
These are exactly all five conditions in (AFE.25).
\end{proof}
% END AFE17

\paragraph{The literal ordinary-point and normalized-representation data.}
Let \(w_0=1,w_1=0,w_2=-1\), and let
\(J\) be the \(3\)-by-\(3\) antidiagonal matrix
whose antidiagonal entries are \(1\).
The group \(\texttt{GG}\,3\,A\) is
\(\mathrm{GL}_3(A)\rtimes\{\pm1\}\), where the sign
\(-1\) acts by
\[
                         X\longmapsto J(X^{-1})^tJ.
 \]
For \(r:G\to\texttt{GG}\,3\,A\), put \(R(g)=(r(g)).\mathrm{mat}\).
An original \(\texttt{OrdPt}\) consists of the following
five conditions: its sign is \(\eta\), and
\[
 R(g)\bmod\mathfrak m_A
   =\operatorname{diag}\bigl(\Psi(g)^{w_i}\bmod\mathfrak m_A\bigr);
 \]
its matrix entries are locally constant modulo every
\(\mathfrak m_A^n\); it has the original normalization
\[
 R(\sigma)=1,\quad
 \bigl(R(g_1)-R(\sigma g_1^{-1}\sigma)\bigr)_{01}
 =\bigl(R(g_1)-R(\sigma g_1^{-1}\sigma)\bigr)_{10}=0;
                                                        \tag{AFE.26}
\]
away from \(\ell\) it kills inertia elements on which
both \(\eta\) and the full \(\Psi\) are \(1\);
and at the specified \(D_B\) there is \(h\equiv1
\pmod{\mathfrak m_A}\) making the matrices upper
triangular on \(\ker\eta\), with inertia diagonals
exactly the images of \(\Psi^{w_i}\).
The last is the literal condition \(\texttt{CondT}\).

An \(\texttt{NRep}(\Phi,3,A)\) consists of a matrix
representation \(\rho:\Phi.\Gamma_0\to\mathrm{GL}_3(A)\)
and local characters \(\psi_{v,i}:\Phi.G_v\to A^\times\).
It requires
\[
 \begin{gathered}
 \rho(\Phi.\star g)=J\rho(g)^tJ,\\
 \rho(\Phi.g_0)-\rho(\Phi.\star\Phi.g_0)
                         \ \text{is diagonal},\\
 \prod_{i=0}^{2}
   \bigl(\rho(\Phi.\operatorname{loc}_{0,v}(g_i))
                         -\psi_{v,i}(g_i)1\bigr)=0
                       \quad\text{for every }(g_0,g_1,g_2).
 \end{gathered}                                         \tag{AFE.27}
\]
The product is in the displayed order.
Its condition \(\texttt{IsLiftA}\) means adic continuity
of all these entries and characters, residual diagonal
entries and local characters prescribed by
\(\Phi.\Psi^{w_i}\), and zero residual off-diagonal
entries.  These are the source's exact data, rather
than a semisimplified substitute.

% BEGIN AFE18
\paragraph{AFE18. Transporting the actual triangular characters.}
\label{claim:AFE18}
Let \(r:G\to\texttt{GG}\,3\,A\) have sign
\((r(g)).\mathrm{right}=\eta(g)\), the residual
diagonal condition, adic continuity and
\(\texttt{CondT}\) at subgroups \(D,I\).
Let \(\varphi:H\to G\) be continuous, and suppose
\[
 \eta(\gamma)=1,\quad \eta(\varphi(s))=1,\quad
                         \gamma^{-1}\varphi(s)\gamma\in D
                           \quad(s\in H).
 \]
Then there are characters \(\chi_i:H\to A^\times\)
with the ordered product condition for \(R\varphi\),
adic continuity, residual values
\(\Psi(\varphi(s))^{w_i}\), and
\[
 \chi_i(s)=\alpha(\Psi(\varphi(s))^{w_i})
 \quad\text{if }\gamma^{-1}\varphi(s)\gamma\in I,
                                                        \tag{AFE.28}
\]
where \(\alpha:O\to A\) is the given local coefficient map.
Here \(A\) need only be a commutative local ring;
no flatness or completeness is used.

\begin{proof}
Choose the actual \(h\equiv1\pmod{\mathfrak m_A}\)
from CondT.  Put
\[
 d(s)=\gamma^{-1}\varphi(s)\gamma,\qquad
 T(s)=h^{-1}R(d(s))h,\qquad \chi_i(s)=T(s)_{ii}.
 \]
All relevant signs are \(1\), so \(T\) is a matrix
representation.  Its matrices are upper triangular.
For such matrices the diagonal of a product is the
product of the diagonals: in the sum
\(\sum_j T(s)_{ij}T(t)_{ji}\), terms \(j<i\)
vanish in the first factor and terms \(j>i\)
in the second.  Consequently each \(\chi_i\)
is a homomorphism; evaluating at \(s^{-1}\)
shows its values are units.

Let \(E_j\) be the span of the first \(j\) coordinate
vectors, with \(E_0=0\).
Upper triangular matrices preserve this filtration.
The factor \(T(g_i)-\chi_i(g_i)1\) maps
\(E_{i+1}\) to \(E_i\).
Thus the rightmost factor for \(i=2\) maps \(A^3\)
to \(E_2\), the middle one maps \(E_2\) to \(E_1\),
and the leftmost maps \(E_1\) to \(0\).
Their ordered product is zero.
Since \(\eta(\gamma)=1\),
\[
 R(\varphi(s))=P\,T(s)\,P^{-1},\qquad P=R(\gamma)h.
 \]
Conjugating the product proves the required version
for \(R\varphi\).

Modulo \(\mathfrak m_A\), both \(h\) and \(h^{-1}\)
are \(1\).
The residual condition therefore gives
\(\bar\chi_i(s)=
 \overline{\alpha(\Psi(d(s))^{w_i})}\).
Since \(\Psi\) is a character,
\(\Psi(d(s))=\Psi(\varphi(s))\).
The exact inertia diagonals of CondT give the same
equality before reduction whenever \(d(s)\in I\).
This proves (AFE.28) with the original character,
not only its residue.

Finally
\(T(s)=P^{-1}R(\varphi(s))P\).
After reduction modulo each \(\mathfrak m_A^n\),
this is a fixed matrix expression in the locally
constant entries of \(R\varphi\).
Its diagonal entries are locally constant as well.
This proves precisely the required adic continuity,
without making any topological assumption on
coefficient changes beyond those stated.
\end{proof}
% END AFE18

% BEGIN AFE19
\paragraph{AFE19. The full original ordinary-point-to-\(\texttt{NRep}\) descent.}
\label{claim:AFE19}
Let \(O\) be a commutative local ring with residue
characteristic \(\ell\).
Take an original
\[
 k_0:\operatorname{GalLink}(\eta,\Psi,\sigma,g_1,B,M,L,\Phi),
 \quad \Phi.\texttt{Ok}(C_0),\quad \ell\mid M,\quad\sigma^2=1.
 \]
Let \(A\) be any Noetherian local domain, with a local
\(O\)-algebra map \(\alpha\).
For every original \(\texttt{OrdPt}\) \(r\), there is
an actual \(\texttt{NRep}(\Phi,3,A)\) satisfying
\(\texttt{IsLiftA}\) and
\[
 \begin{aligned}
 \rho(k_0.p_0(s))&=R(s) &&(s\in G_L),\\
 \psi_{v,i}(t)&=\alpha(\Psi(k_0.j_v(t))^{w_i})
                                      &&(t\in\Phi.I_v).
 \end{aligned}                                         \tag{AFE.29}
\]
The inclusion \(G_L\to G\) is understood in the first
line.  This is the full original
\(\texttt{exists\_nrep3}\), with no extra arithmetic
or representation-theoretic premise.

\begin{proof}
Put \(U=\operatorname{upl}(L)\) and \(p=k_0.p:U\to\Phi.\Gamma\).
The subgroup \(U\) is closed in the profinite \(G\):
it is an intersection of stabilizers of algebraic
elements, and each such stabilizer is open and closed
because the orbit is finite.
Thus \(U\) is compact.
The exact GalLink kernel is the closure of the normal
subgroup generated by those \(t\in U\) for which
\[
 \exists B'\nmid\ell:\quad
 t\in I_{B'},\quad \eta(t)=1,\quad\overline\Psi(t)=1.
                                                        \tag{AFE.30}
\]
Its Teichmuller field gives \(\Psi(t)=1\) for each
generator in (AFE.30), and the away field of OrdPt
then gives \(r(t)=1\).
Hence \(r\) kills their normal subgroup.

We justify passage to its closure.
For the Noetherian local ring \(A\),
\(\bigcap_n\mathfrak m_A^n=0\).
One standard algebraic proof applies Artin--Rees
to the ideal \(I=\bigcap_n\mathfrak m_A^n\subset A\):
for a sufficiently large integer \(c\),
\(I\cap\mathfrak m_A^{c+1}\subset\mathfrak m_A I\).
Since \(I\) lies in every power, \(I=\mathfrak m_A I\);
it is finitely generated, so Nakayama gives \(I=0\).
Therefore
\[
 \{u\in U:R(u)=1\}=
 \bigcap_n\{u\in U:R(u)\equiv1\pmod{\mathfrak m_A^n}\}
 \]
is closed: each set on the right is a fiber of a
locally constant map.  It contains the normal
subgroup just considered, hence its closure.
For \(u\in\ker p\), GalLink's sign relation also gives
\(\eta(u)=\Phi.\delta(pu)=1\).
Both the matrix and the sign of \(r(u)\) are therefore
trivial.  Thus \(\ker p\subset\ker(r|_U)\).

Surjectivity of \(p\) gives a unique homomorphism
\[
                  r':\Phi.\Gamma\to\texttt{GG}\,3\,A,
                         \qquad r'(pu)=r(u).
 \]
A continuous surjection from compact \(U\) to the
Hausdorff profinite \(\Phi.\Gamma\) is a quotient map:
it is closed, so a subset of the target is open
exactly when its inverse image is open.
Applying this to the fibers of each matrix reduction
shows that \(r'\) is adically continuous.
The sign equality \((r'(x)).\mathrm{right}=\Phi.\delta(x)\)
follows by surjectivity from GalLink's sign relation.
Hence
\[
             \rho(g)=(r'(g)).\mathrm{mat}
                         \quad(g\in\Phi.\Gamma_0)
 \]
is a matrix representation.

The GalLink fields at \(\sigma,g_1\) give
\(r'(\Phi.c)=r(\sigma)\) and
\(r'(\Phi.g_0)=r(g_1)\).
Ok gives \(\Phi.\delta(\Phi.c)=-1\) and \(\Phi.c^2=1\);
normalization gives \(R(\sigma)=1\).
In the displayed semidirect product,
\[
 (1,-1)(X,1)^{-1}(1,-1)=(JX^tJ,1).
 \]
Since \(\Phi.\star g=\Phi.c\,g^{-1}\Phi.c\),
this proves the polarization in (AFE.27).
For any matrix \(X=(x_{ij})\),
\[
 X-JX^tJ=
 \begin{pmatrix}
 x_{00}-x_{22}&x_{01}-x_{12}&0\\
 x_{10}-x_{21}&0&x_{12}-x_{01}\\
 0&x_{21}-x_{10}&x_{22}-x_{00}
 \end{pmatrix}.
 \]
Its two entries \(01,10\) control all off-diagonal
entries.  At \(X=R(g_1)\), (AFE.26) makes both zero,
so the normalization field of NRep follows.

For each \(v\), the original
\(\texttt{pl\_conj}\) field provides \(\gamma_v\)
with \(\eta(\gamma_v)=1\) and
\(k_0.\mathrm{pl}_v=\gamma_v B\).
Set \(\varphi=k_0.j_v\), viewed as a map to \(G\).
The fields \(\texttt{j\_range}\) and \(\texttt{j\_ine}\)
identify its image and inertia with the actual
decomposition and inertia at this place.
Conjugation by \(\gamma_v^{-1}\) therefore carries
them into \(D_B,I_B\).
Furthermore GalLink's sign and localization equalities
give
\[
           \eta(k_0.j_v(t))
             =\Phi.\delta(\Phi.\operatorname{loc}_v(t))=1,
 \]
since this localization has its specified lift to
\(\Phi.\Gamma_0\).
Apply AFE18 to obtain \(\psi_{v,i}\).
GalLink's localization equality and the definition of
\(r'\) give
\[
          \rho(\Phi.\operatorname{loc}_{0,v}(t))
                           =R(k_0.j_v(t)).
 \]
Thus AFE18 gives the ordered product field in
(AFE.27), continuity, and the second line of (AFE.29).

It remains to check the literal residual fields.
The local coefficient map induces
\(k\to k_A=A/\mathfrak m_A\), characterized by
\(\operatorname{red}_A\alpha=
 (k\to k_A)\operatorname{red}_O\).
For a lift \(u\in U\) of \(g\in\Phi.\Gamma_0\),
OrdPt gives the residual diagonal
\(\alpha(\Psi(u)^{w_i})\), and GalLink gives
\(\Phi.\Psi(pu)=\overline\Psi(u)\).
These are exactly the prescribed diagonal entries;
the off-diagonal entries vanish by the same OrdPt
condition.  AFE18 gives the corresponding local
character identities.  All weights are the literal
\(\texttt{w}\,3\,i=1-i=w_i\).
Together with the adic continuity already proved,
these are all fields of IsLiftA.
Finally \(k_0.p_0\) is the restriction of \(p\)
to \(G_L\), so \(r'(pu)=r(u)\) gives the first
line of (AFE.29).
\end{proof}
% END AFE19

% BEGIN AFE20
\paragraph{AFE20. The full original \(\texttt{FinPt}\Rightarrow\texttt{FinEnd}\) transport.}
\label{claim:AFE20}
At the exact GalLink, Ok, residue-characteristic,
\(\ell\mid M\), and \(\sigma^2=1\) premises of
AFE19, the original \(\texttt{FinPt}(O,\Phi,3)\)
implies
\[
                    \texttt{FinEnd}(\ell,\eta,\Psi,\sigma,g_1,B).
                                                        \tag{AFE.31}
\]
The finite-index subgroup can be chosen as
\(G_L\) before every coefficient ring and every
ordinary point.

\begin{proof}
Choose \(U_0=G_L\subset G\) once.
Its index is \([L:\mathbb Q]\), a positive finite
integer.  This choice depends only on the given
frame and field.
Now take any original \(\texttt{IsCD}(O,A)\)
coefficient algebra and its local map
\(\alpha:O\to A\), with
\(\alpha(\mathfrak m_O)=0\), and any ordinary point
\(r:G\to\texttt{GG}\,3\,A\).
Here IsCD means a complete DVR domain \(A\), a
local coefficient map, and finite residue field;
in particular \(A\) is a Noetherian local domain.
There is a unique ring map \(f:k\to A\) with
\[
                         f\operatorname{red}_O=\alpha.
 \]
Give \(A\) this \(k\)-algebra structure.
It forms the required scalar tower \(O\to k\to A\).

AFE19 supplies \(\rho\) with IsLiftA and (AFE.29).
Suppose \(t\in\Phi.I_v\) and
\(\Phi.\Psi(\Phi.\operatorname{loc}_v(t))=1\).
GalLink's character relation makes this
\(\overline\Psi(k_0.j_v(t))=1\).
Because \(\alpha\) factors through \(k\), the exact
inertia formula in (AFE.29) gives
\[
                    \psi_{v,i}(t)
                       =f\bigl(\overline\Psi(k_0.j_v(t))^{w_i}\bigr)
                       =1.
 \]
This is precisely the extra inertia condition in
the original FinPt, including the negative weight.
FinPt therefore gives, for every \(g\in\Phi.\Gamma_0\),
\[
       \operatorname{charpoly}(\rho(g))
         =\prod_{i=0}^{2}
                 \left(X-f(\Phi.\Psi(g)^{w_i})\right).
                                                        \tag{AFE.32}
\]

Take any \(s\in U_0\), retaining the original
\(\eta(s)=1\) condition of FinEnd.
Represent it as an element of \(G_L\) and apply
(AFE.32) to \(k_0.p_0(s)\).
By (AFE.29) the matrix on the left is \(R(s)\).
The character relation and \(f\operatorname{red}_O=\alpha\)
identify the right side with
\[
                  \prod_{i=0}^{2}(X-\alpha(\Psi(s)^{w_i})).
 \]
The matrix of the original residual reference
representation \(r_0(\eta,\Psi)(s)\) is
\(\operatorname{diag}(\Psi(s),1,\Psi(s)^{-1})\).
After applying \(\alpha\), its characteristic
polynomial is exactly this product.
This proves the original FinEnd equality.
Since \(U_0\) was chosen before \(A,r,s\), its
quantifier order is the required one.
The proof has used the supplied FinPt and has not
asserted its separate patching-theoretic producer.
\end{proof}
% END AFE20

% BEGIN AFE21
\paragraph{AFE21. The common marked frame now carries the original finite-end implication.}
\label{claim:AFE21}
Under \(P\), each original Dat input, and each original
HostData input with its specified \(\sigma,B\), admits
the same output \(L,g_1,\Phi,C_0\) as in AFE15 or
AFE16, together with its full original Marks and
the implication
\[
          \texttt{FinPt}(c.O,\Phi,3)
              \ \Longrightarrow\
          \texttt{FinEnd}(\ell,\eta,\Psi,\sigma,g_1,B).
                                                        \tag{AFE.33}
\]
The \(\eta,\sigma,M,H_b\) in the two cases retain
their respective original meanings.

\begin{proof}
Choose the output only once, by AFE15 or AFE16.
AFE17 supplies Marks on its \(g_1\), including
\(\sigma^2=1\).
The actual modulus is divisible by \(\ell\) in
both cases.  The GalLink and all Ok fields are
those already chosen, and AFE1 supplies the
residue characteristic.
Apply AFE20 to this same data.
Its subgroup is exactly \(G_L\), independent of
the later coefficient algebra and ordinary point.
This proves (AFE.33) while preserving \(Q\) before
\(L,w\), the same \(g_1\), and the same frame.
No simultaneous choice involving an automorphic
representation has been added.
\end{proof}
% END AFE21

\paragraph{Boundary after these joins.}
The original cohomological premise \(P\) remains in
AFE15--16.  The resulting SeedClause supplies the
specified arithmetic characters; it does not supply
the separate rank-two and rank-three automorphic
seed representations.  Those original producers,
and the later \(\texttt{FinPt}\), \(\texttt{AutPt}\),
or their distinct Alt versions, remain explicit
inputs at their later applications.  The two Gen
entries, the full HasFrame conclusions at \(P\),
the actual auxiliary Host-frame witness statements,
the normalized-representation transport, and
\(\texttt{FinPt}\Rightarrow\texttt{FinEnd}\)
are complete at their original premises.
The full common \(\texttt{GdEnd}\) and
\(\texttt{HostEnd}\) still require their remaining
finite-point and automorphic suppliers.
